\begin{table}[H]
\centering
\captionsetup{justification=centering}
\caption{Correlação estimada entre os distúrbios das inclinações - modelo multivariado - total de ocupados}
\label{tab:matriz_correlacao_ocup}

\renewcommand{\arraystretch}{0.8}

\resizebox{\textwidth}{!}{%
\begin{tabular}{@{}p{2.8cm}c*{8}{c}@{}}
\toprule
\textbf{Estrato Geográfico} & \(\hat{\sigma}_R^2\) & \textbf{1} & \textbf{2} & \textbf{3} & \textbf{4} & \textbf{5} & \textbf{6} & \textbf{7} & \textbf{8} \\
\midrule
01 - Belo Horizonte & 0,8061 & 1 &  &  &  &  &  &  &  \\
\raisebox{0em}{\parbox[t]{4.3cm}{02 - Entorno e Colar \\[-2pt] Metropolitano de BH}} & 7,0885 & 0,2058 & 1 &  &  &  &  &  &  \\
03 - Sul de Minas & 5,0512 & 0,2438 & 0,2887 & 1 &  &  &  &  &  \\
04 - Triângulo Mineiro & 8,7550 & 0,1852 & 0,2193 & 0,2641 & 1 &  &  &  &  \\
05 - Zona da Mata & 11,6196 & 0,1607 & 0,1904 & 0,2293 & 0,2415 & 1 &  &  &  \\
06 - Norte de Minas & 3,8825 & 0,2781 & 0,3294 & 0,3966 & 0,4179 & 0,4421 & 1 &  &  \\
07 - Vale do Rio Doce & 1,8018 & 0,4082 & 0,4835 & 0,5822 & 0,6134 & 0,6490 & 0,8872 & 1 &  \\
08 - Central & 3,3652 & 0,2987 & 0,3538 & 0,4260 & 0,4488 & 0,4749 & 0,6492 & 0,7363 & 1 \\
\bottomrule
\end{tabular}%
}

\fonte{Elaboração própria, com base nos dados da PNAD Contínua. Nota: a matriz de covariância dos distúrbios das inclinações é parametrizada por um fator de Cholesky (\(\Sigma_R = LL'\)), o que garante que a matriz estimada seja positiva semidefinida por construção. Autovalores estimados: 19,795; 7,915; 6,116; 4,066; 2,485; 1,281; 0,710 e 1 valor inferior a \(10^{-3}\). Posto numérico de 8 em 8; posto \textbf{efetivo} de 7, contando os autovalores que retêm ao menos 1\% do maior. As correlações devem, portanto, ser lidas em conjunto, como indicação de um número reduzido de tendências comuns, e não isoladamente.}
\end{table}
